\section{Equa��o 1}

Seja a seguinte equa��o diferencial ordin�ria de segunda ordem, sendo $u = u(x)$, e seja
$u'$ e $u''$ as primeira e segunda derivadas da fun��o $u(x)$ em rela��o � x:
%
%
\[u'' - xu' +\lambda u = 0\]
%
%
Nesta se��o demonstraremos que a fun��o, $v(x) = xu - u'$, satisfaz a seguinte equa��o
diferencial:
%
%
\[v'' - xv' + (1 + \lambda) v = 0\]
%
%
Temos que:
%
%
\begin{eqnarray}
\nonumber  v'  &=& (xu-u')' = u + xu' - u''\\
\nonumber  v'' &=& (xu-u')'' = (u + xu' - u'')' = 2u' + xu'' - u'''
\end{eqnarray}
%
%
Inserindo as rela��es acima, na equa��o diferencial, temos que:
%
%
\[v'' - xv' + (1 + \lambda) v  = (xu-u')' = 2u' + xu'' - u''' - x(u + xu' -
u'') + (1 + \lambda)(xu - u')\]
%
%
Simplificando as rela��es acima, encontramos que:
%
%
\[v'' - xv' + (1 + \lambda)v = u' + xu'' - u''' - x^2u' + xu'' + \lambda xu - \lambda u'\]
%
%
Reorganizando de maneira apropriada, encontramos que:
%
%
\[v'' - xv' + (1 + \lambda) v = u' + xu'' - u''' + x(u'' - xu' + \lambda u) - \lambda u'\]
%
%
Lembrando que, $u'' - xu' + \lambda u = 0$, e inserindo na rela��o acima, temos que:
%
%
\[v'' - xv' + (1 + \lambda) v = u' + xu'' - (u''' + \lambda u')\]
%
%
Notando que, $u' + xu'' = (xu')'$, e que $u''' + \lambda u' = (u'' + \lambda u)'$,
encontramos que:
%
%
\[v'' - xv' + (1 + \lambda) v = (xu')' - (u'' + \lambda u)'\]
%
%
Aplicando a linearidade da derivada, temos finalmente que:
%
%
\[v'' - xv' + (1 + \lambda) v = - (- xu' + u'' + \lambda u)' = 0\]

